\documentstyle[% 


righttag% 


,11pt% 


,amssymb% 


,russian%               remove in Eng. 


,izvru%                 Set izven% in Eng. 


]{amsart} 


 


%       Math definitions 


 


\newcommand{\sign}{\operatorname{sign}} 


\newcommand{\tts}[1]{} 


 


\theoremstyle{plain} 


\theorembodyfont{\it} 


\newtheorem{lemnonum}{\hskip\parindent Лемма} 


\renewcommand{\thelemnonum}{} 


 


\theoremstyle{definition} 


\theorembodyfont{\rm} 


\newtheorem{cornonum}{\hskip\parindent Следствие} 


\newtheorem{notenonum}{\hskip\parindent Замечание} 


\renewcommand{\thecornonum}{} 


\renewcommand{\thenotenonum}{} 


 


%\hoffset=-15mm%                  For Eng. 


  


\setcounter{page}{32} 


\god{1999} 


\nomer{4~(443)} %                        Remove in Eng. 


%\nomer{Vol.\,43, No.\,4}%               Set in Eng. 


%\str{\pageref{begin}--\pageref{end}}%  Set in Eng. 


\udk{517.518} 


\author{\it Н.Ш.\,Загиров} 


% NB:   ^^ remove in Eng. 


\title{Полиномиальные ужи по подсистемам\\ 


алгебраических степеней} 


 


\date{} 


% \pagestyle{myheadings}%         Set in Eng. 


\pagestyle{plain} %              Remove in Eng. 


\begin{document} 


% \thispagestyle{plain}%          Set in Eng. 


\maketitle 


\label{begin} 


 


 


{\bf 1. Обозначения. Постановка задачи.} Полиномиальные ужи, 


введенные в [1], получили дальнейшее 


развитие в работах В.К.\,Дзядыка (см., напр., [2]). Их обобщение 


на каналы с разрывными границами построено в [3]. 


 


Пусть на конечном отрезке $[a,b]$ задана некоторая линейно 


независимая система функций $\varphi_0,\dotsc,\varphi_n$. Через $L$ 


обозначим множество полиномов по этой системе. Для двух фиксированных 


функций $g$, $G$ таких, что $g(x)<G(x)$, когда $x$ принадлежит 


$[a,b]$, положим $L(g,G)=\{P\in L:g(x)\le P(x)\le G(x),\ x\in[a,b]\}$. 


 


Напомним, что полином $Z\in L(g,G)$ называется ужом, 


порожденным парой $(g,G)$, или ужом пары $(g,G)$, если на отрезке 


$[a,b]$ имеется такой набор из $m=n+1$ точек $x_1<x_2<\dotsb<x_m$, 


что с ростом $k$ от 1 до $m$ точки $M_k=(x_k,Z(x_k))$ попеременно 


попадают то на график $\Gamma(g)$ функции $g$, то на график $\Gamma(G)$ 


функции $G$ (начиная с любого из них). Если же $m>n+1$, то $Z$ 


называют избыточным ужом [3]. С.\,Карлин доказал утверждение: 


если $g$, $G$ --- непрерывные на отрезке $[a,b]$ функции, для которых 


существует такой полином $l$ из $L$, что $g(x)<l(x)<G(x)$ на 


$[a,b]$ и система $\varphi_0,\dotsc,\varphi_n$ 


является чебышевской, то имеются 


ровно два ужа, порожденные парой $(g,G)$. 


 


Этот результат обобщен Е.П.\,Долженко и Е.А.\,Севастьяновым [2] 


на разрывные функции $g$, $G$, ими же установлены некоторые 


новые свойства ужей, а также их связь с аппроксимациями 


со знакочувствительным весом [4]--[6]. 


 


В данной работе изучаются вопросы существования и 


количества ужей, когда система $\varphi_0,\dotsc,\varphi_n$ не является 


чебышевской; 


рассматривается одна экстремальная задача, связанная 


с ужами. Установленные результаты конкретизируются на 


полиномах по некоторым подсистемам алгебраических степеней. 


\medskip 


 


{\bf 2.} Будем считать все рассматриваемые ниже функции непрерывными 


на отрезке $[a,b]$. С парой $(g,G)$ свяжем две новые 


функции: $g_0(x)=(g(x)+G(x))/2$, $h(x)=(G(x)-g(x))/2$. 


Так как $G(x)>g(x)$, то $h(x)>0$ на отрезке $[a,b]$. 


 


Введем функционал 


$$ 


N(f)=N(f,g,G)=\max\{|f(x)-g_0(x)|/h(x):x\in[a,b]\}. 


$$ 


Известно [7], [8], что $N(f)$ --- непрерывный выпуклый неотрицательный 


функционал и 


\begin{equation} 


L(g,G)=\{l\in L:N(l)\le 1\}. 


\end{equation} 


Величину $N(f_1-f_2)$ называем $(g,G)$-расстоянием между 


$f_1$ и $f_2$. Пусть 


$$ 


E_n(f)=\inf_{l\in L}\max\{|f(x)-l(x)|/h(x):x\in[a,b]\}. 


$$ 


Напомним, что величина $E_n(f)$ называется наименьшим уклонением 


пространства $L$ от функции $f$, а $l^*\in L$ такой, что 


$$ 


\max\{|f(x)-l^*(x)|/h(x):x\in[a,b]\}=E_n(f)


$$ 


называется полиномом (элементом) наилучшего приближения $f$ на 


$L$. Очевидно, для положительных на отрезке $[a,b]$ функций $h$ 


сохраняются теоремы существования, единственности, характеристики 


элемента наилучшего приближения, известные в случае 


$h=1$ [3]. Для общего случая, когда 


$F_n(f)=F_n(f;g,G)=\inf\{N(f-l):l\in L\}$, 


сохраним ту же терминологию. 


 


Отметим очевидное равенство $F_n(f)=E_n(f-g_0)$, 


из которого следует, что $l^*$ из $L$ будет полиномом наилучшего 


приближения $f$ в смысле $(g,G)$-расстояния тогда и только тогда, 


когда он будет полиномом наилучшего приближения для 


функции $f-g_0$ в обычном смысле. Это замечание доказывает, 


что упомянутые теоремы распространяются и на общий случай. 


 


В этой работе ужи будут построены как решение некоторой 


экстремальной задачи. Именно, фиксируем одну из функций системы 


$\varphi_0,\dotsc,\varphi_n$, скажем, $\varphi_n$. Полином $l$ из $L$ 


будем записывать 


в виде $l=c_n(l)\varphi_n+q$, где $q$ принадлежит множеству $V$ 


полиномов, порожденных системой $\varphi_0,\dotsc,\varphi_{n-1}$. 


 


Положим 


\begin{equation} 


M=\sup\{c_n(l):l\in L(g,G)\},\quad 


m=\inf\{c_n(l):l\in L(g,G)\}.\tag{\text{з}1} 


\end{equation} 


Так как $L(g,G)$ есть выпуклый компакт, то обе задачи имеют 


решения, причем, если $l^*$ --- решение какой-нибудь из них, то 


$N(l^*)=1$. Чтобы oхарактеризовать элемент $l^*$, введем новые 


обозначения. Для фиксированного числа $t$ положим 


\begin{gather*} 


F(t)=\inf\{N(t\varphi_n-q):q\in V\},\\ 


F=F(g,G)=\inf\{F(t):-\infty<t<\infty\}. 


\end{gather*} 


Очевидно, 


\begin{gather*} 


F(t)=\inf\{N(l):c_n(l)=t\},\\ 


F=F(g,G)=\inf\{N(l):l\in L(g,G)\}. 


\end{gather*} 


Величину $F(g,G)$ будем называть расстоянием от пространства $L$ 


до канала $(g,G)$. Заметим, что $F(g,G)=0$ тогда и только 


тогда, когда ``середина'' канала $g_0$ принадлежит $L$. 


 


Нам понадобится 


 


\begin{lemnonum} Функция $F(t)$ является непрерывной выпуклой на 


всей числовой прямой, причем 


\begin{equation} 


\lim_{|t|\to\infty}F(t)=+\infty. 


\end{equation} 


 


Множество $L(g,G)$ не пусто тогда и только тогда, когда $F(g,G)\le 1$. 


\end{lemnonum} 


 


\begin{pf} Последняя часть утверждения фактически 


следует из равенства (1). Выпуклость $F$ следует из выпуклости 


функционала $N$:


$$ 


F((t_1+t_2)/2)\le N(\varphi_n(t_1+t_2)/2- 


(q_1+q_2)/2) \le \tfrac{1}{2}(N(t_1\varphi_n-q_1)+ 


N(t_2\varphi_n-q_2))=\tfrac{1}{2}(F(t_1)+F(t_2)). 


$$ 


Здесь $q_i\in V$ --- полином наилучшего приближения функции 


$t_i\varphi_n$, 


$i=1,2$. Если $g_0\notin L$, то равенство (2) непосредственно 


получается 


из неравенства ($t\ne 0$) 


\begin{equation} 


F(t)\ge |t|e_n-K(g,G), 


\end{equation} 


где $e_n=\inf\{N(\varphi_n-q):q\in V\}$, 


$K(g,G)=\max\{|g_0(x)|/h(x):x\in[a,b]\}$. 


Так как система функций $\varphi_0,\dotsc,\varphi_n$ линейно 


независимая и 


$g_0\notin L$, то $e_n>0$. 


Неравенство (3) получается из следующих двух неравенств: 


\begin{gather*} 


(|t\varphi_n(x)-q^*(x)|-|g_0(x)|)/h(x)\le N(t\varphi_n-q^*),


\quad x\in [a,b]; \\ 


\sup\{(|t\varphi_n(x)-q^*(x)|-|g_0(x)|)/ 


h(x):x\in[a,b]\}\ge \\ 


%


\ge\sup\{(|t\varphi_n(x)-q^*(x)|)/h(x): 


x\in[a,b]\}-\sup\{|g_0(x)|/h(x):x\in[a,b]\}. 


\end{gather*} 


Здесь $q^*$ --- полином наилучшего приближения функции 


$t\varphi_n$. Если $g_0\in L$, то достаточно считать 


$g_0(x)=t_0\varphi_n(x)$ и $t\ne t_0$. 


\end{pf} 


 


\begin{cornonum} Для любой пары $(g,G)$ существуют такие числа 


$t_1$, $t_2$, что на луче $(-\infty,t_1)$ функция $F(t)$ строго убывает, на 


луче $(t_2,\infty)$ 


строго возрастает, а на отрезке $[t_1,t_2]$ она равна 


$F(g,G)$. Если система функций $\varphi_0,\dotsc,\varphi_n$ является 


чебышевской, 


то $t_1=t_2$. 


\end{cornonum} 


 


Докажем теперь теорему о характеристике решения задач (з1). 


 


\begin{thm} Пусть $L(g,G)\ne 0$. Тогда полином $l^*$ из $L$ будет 


решением 


первой из задач {\em (з1)}, т.\,е. 


\begin{equation} 


c_n(l)\to\sup,\ \;l\in L(g,G),\tag{$*$} 


\end{equation} 


тогда и только тогда, когда 


$q^*=M\varphi_n-l^*$ $(M$ из {\em (з1))} 


является полиномом наилучшего приближения функции $M\varphi_n$ в $V$ 


в смысле $(g,G)$-расстояния. 


 


Аналогичное утверждение верно и для второй задачи {\em (з1)}. 


\end{thm} 


 


\begin{pf} Так как $L(g,G)\ne 0$, то $F(g,G)\le 1$. Пусть 


$F(g,G)=1$. Для любого $t<t_1$ и $t>t_2$ функция $F(t)>1$. Поэтому 


для любого полинома $l$ вида $l=t\varphi_n+q$ число $N(l)>1$, а значит 


(см. равенство (1)), $l$ не принадлежит $L(g,G)$ и $M=t_2$, $m=t_1$. 


Если $l^*=M\varphi_n+q^*$ --- решение задачи ($*$), то $l^*$ 


принадлежит $L(g,G)$ и т.\,к.~$F(g,G)=1$, то $F(M)=F(t_2)=1$. Из 


определения 


$F(t)$ следует, что $q^*$ --- полином наилучшего приближения для 


$M\varphi_n$ 


в смысле $(g,G)$-расстояния. 


 


Допустим, что $q^*=M\varphi_n-l^*$ из $V$ есть полином наилучшего 


приближения для $M\varphi_n$. Тогда $l^*=M\varphi_n-q^*$ и 


$F(M)=F(t_2)=1$. Значит, $l^*$ принадлежит $L(g,G)$ 


и является решением задачи ($*$). 


Аналогично разбирается вторая задача. 


 


Рассмотрим случай, когда $F(g,G)<1$. Из леммы следует, 


что существуют ровно два таких числа $a_1<t_1\le t_2<a_2$, что 


$F(a_1)=F(a_2)=1$. Поскольку для $t<a_1$ и $t>a_2$ имеем $F(t)>1$ и 


полином $l=t\varphi_n+q$ 


не принадлежит $L(g,G)$, то $M=a_2$, $m=a_1$. Теперь 


из определения $F(a_1)$ и $F(a_2)$ следует утверждение теоремы.


\end{pf} 


 


Фактически, как следствие доказанной теоремы и теоремы 


Чебышева о характеристике полинома наилучшего приближения 


получается


 


\begin{thm} Пусть $F(g,G)=1$. Тогда {\em a)}~если система функций 


$\varphi_0,\dotsc,\varphi_n$ является чебышевской, то канал $(g,G)$ 


содержит 


единственный полином, который для него является избыточным 


ужом $(t_1=t_2);$ {\em б)}~если же подсистема 


$\varphi_0,\dotsc,\varphi_{n-1}$ является чебышевской, 


то канал $(g,G)$ имеет единственный уж, когда функция 


$g_0$ имеет единственный полином наилучшего приближения в 


$L$. Канал имеет бесконечно много ужей, если $g_0$ имеет хотя 


два полинома наилучшего приближения в $L$. 


 


Пусть $F(g,G)<1$. Тогда, если система 


$\varphi_0,\dotsc,\varphi_{n-1}$ является 


чебышевской, то пара $(g,G)$ порождает по меньшей мере два ужа по 


системе $\varphi_0,\dotsc,\varphi_n$. 


\end{thm} 


 


\begin{pf} Так как $F(t)=F_n(t\varphi_n)=E_n(t\varphi_n-g_0)$, 


то в доказательстве нуждается лишь та часть теоремы, где 


$F(g,G)=1$, а система $\varphi_0,\dotsc,\varphi_{n-1}$ является 


чебышевской. Пусть 


$g_0$ имеет в $L$ единственный полином $l^*$ наилучшего приближения. 


Так как $F(g,G)=1$, то $E_n(g_0)=1$. Если считать, что 


$l^*=t^*\varphi_n+q^*$, 


то $q^*$ будет полиномом наилучшего приближения для функции 


$g_0-t^*\varphi_n$ 


по чебышевской системе $\varphi_0,\dotsc,\varphi_{n-1}$. 


Следовательно, по 


теореме Чебышева [2] существуют точки $x_1<\dotsb<x_m$ 


($m\ge n+1$), в которых $g_0(x_i)-l^*(x_i)=\varepsilon(-1)^i h(x_i)$, 


$i=1,\dotsc,m$, 


где $\varepsilon$ равен либо плюс, либо минус единице. Возвращаясь 


от $g_0$ и $h$ к $g$ и $G$, можно убедиться, что $l^*$ является 


ужом пары $(g,G)$. 


 


В силу единственности полинома наилучшего приближения 


для каждого $l\in L$, $l\ne l^*$, будем иметь $N(l)>1$ (если 


$N(l)=1$, 


то $l$ тоже был бы полиномом наилучшего приближения $g_0$ в $L$), 


т.\,е.~$l\in L(g,G)$. Следовательно, канал $(g,G)$ 


содержит единственный полином $l^*$. 


 


Пусть теперь $g_0$ допускает хотя бы два полинома 


$l_1=a_1\varphi_n+q_1$, $l_2=a_2\varphi_n+q_2$ наилучшего приближения. 


Так как система 


$\varphi_0,\dotsc,\varphi_{n-1}$ является чебышевской, то $a_1\ne a_2$. 


Как обычно, можно убедиться, 


что тогда любой полином $l_z=z l_1+(1-z)l_2$, 


$z\in[0,1]$, 


также является полиномом наилучшего приближения для $g_0$ и 


$E_n(g_0)=1$. 


Это же означает, что для каждого $z\in[0,1]$ 


полином $z q_1+(1-z)q_2$ является наилучшим в чебышевском 


пространстве $V$ для функции $g_0=z a_1+(1-z)a_2$. Следовательно


(см. предыдущий абзац), каждый из полиномов $l_z$ является ужом 


пары $(g,G)$. 


\end{pf} 


 


\begin{cor} Если для пары $(g,G)$ выполняется неравенство $F(g,G)<1$, 


а система $\varphi_0,\dotsc,\varphi_{n-1}$ является чебышевской, то 


решения задач (з1) являются ужами, порожденными парой $(g,G)$. 


Если же и система $\varphi_0,\dotsc,\varphi_{n-1},\varphi_n$ является 


чебышевской, то 


два ужа, порожденные парой $(g,G)$, и только они будут 


решениями задач (з1). 


\end{cor} 


 


\begin{note} Вторая часть теоремы 2 показывает: чтобы канал $(g,G)$ 


имел по меньшей мере два ужа, необязательно, чтобы вся система 


$\varphi_0,\dotsc,\varphi_n$ была чебышевской. Для этого достаточно, 


чтобы 


существовала хотя бы одна чебышевская подсистема из $n$ 


функций (для $F(g,G)<1$). (Ср. с теоремой С.\,Карлина.) 


\end{note} 


 


\begin{note} Простые примеры ($\varphi_0(x)=x$, $\varphi_1(x)=x^2-1$, 


$\varphi_2(x)=x^2-4$, $n=2$) 


показывают, что чебышевская система $\varphi_0,\dotsc,\varphi_n$ 


может и не содержать чебышевских подсистем из $n$ функций. 


\end{note} 


 


\begin{note} Случай a) первой части теоремы 2 впервые отмечен 


Е.П.\,Долженко и\linebreak Е.А.\,Севастьяновым. 


\end{note} 


 


Из теоремы 2 и теоремы С.\,Карлина получаем 


 


\begin{cor} Если для пары $(g,G)$ выполняется условие


$F(g,G)<1$, а система $\varphi_0,\dotsc,\varphi_n$ или какая-нибудь ее 


подсистема 


из $n$ функций является чебышевской, то пара $(g,G)$ порождает 


по меньшей мере два ужа. 


\end{cor} 


 


Если система $\varphi_0,\dotsc,\varphi_n$ не является чебышевской, то, 


вообще говоря,  


существует пара $(g,G)$, которая порождает по 


этой системе бесконечно много ужей. Так, например, пусть для 


системы $\varphi_0,\dotsc,\varphi_n$ выполняется условие (i): для 


некоторого 


набора точек $x_0<\dotsb<x_n$ отрезка $\Delta=[a,b]$ 


$$ 


\begin{vmatrix} 


\varphi_0(x_0) & \dotsc &\varphi_0(x_n)\\ 


\hdotsfor[1.5]{3} \\


\varphi_n(x_0)&\dotsc &\varphi_n(x_n) 


\end{vmatrix} 


=0, 


$$ 


определители 


$$ 


\Delta_i= 


\begin{vmatrix} 


\varphi_0(x_0) &\dotsc &\varphi_0(x_{i-1}) &


\varphi_0(x_{i+1}) &\dotsc &\varphi_0(x_n)\\ 


\hdotsfor[1.5]{6} \\


\varphi_{n-1}(x_0)&\dotsc &\varphi_{n-1}(x_{i-1}) &


\varphi_{n-1}(x_{i+1})&\dotsc &\varphi_{n-1}(x_n) 


\end{vmatrix} 


$$ 


имеют один и тот же знак. 


 


Можно показать, что при условии (i) 


\begin{itemize} 


\item[1)] существуют такие числа $b_i$ ($i=0,\dotsc,n$), что 


$b_i b_{i+1}<0$ для $i=0,\dotsc,n-1$, и 


$\sum\limits^n_{i=0}b_i P(x_i)=0$  


для любого $P\in L_n$; 


\item[2)] существует полином $P_0\in L_n$ с нулями в точках 


$x_0,x_1,\dotsc,x_n$ 


и такой, что $\max\limits_{x\in[a,b]}|P_0(x)|=1$. 


\end{itemize} 


 


Положим, что $f$ есть непрерывная на отрезке $[a,b]$ функция с 


условиями $f(x_i)=\sign b_i$, 


$\max\{|f(x)|:x\in[a,b]\}=1$. Пусть 


$g_0(x)=f(x)(1-|P_0(x)|)$, $G(x)=g_0(x)+1$, $g(x)=g_0(x)-1$. 


Тогда при $|\varepsilon|\le 1$ любой полином вида 


$\varepsilon P_0(x)$ образует уж для 


пары $(g,G)$. 


 


В работе [9] доказано, что множество систем, удовлетворяющих 


условию (i), непусто. 


 


Если же ни система $\varphi_0,\dotsc,\varphi_n$ и никакая ее 


подсистема из 


$n$ функций не будет чебышевской, то существуют каналы без 


ужей. Например, канал $(g,G)$ с 


\begin{gather*} 


G(x)=g_1(x)+1,\quad g(x)=g_1(x)-5/4, \\ 


g_1(x)= 


\begin{cases} 


(7/4)x+1, & -1\le x\le 0;\\ 


1-(3/4)x, & 0\le x\le 1, 


\end{cases} 


\end{gather*} 


не порождает ужей по системе $x$, $x^3$, $x^5$. 


\medskip 


 


{\bf 3.} В этом пункте изучим свойства подсистем алгебраических 


степеней. Считаем, что 


$$ 


\varphi_i(x)=x^i,\ \; i=0,\dotsc,j-1;\quad 


\varphi_i(x)=x^{i+q},\ \; i=j,\dotsc,n, 


$$ 


где $q\ge 1$, $1\le j\le n+1$. 


 


При $j=n+1$ эта система является чебышевской на любом отрезке. 


Известно также, что она является чебышевской на любом 


отрезке $[a,b]$, когда $ab>0$. Случай $q=1$ рассмотрен автором 


в [9]. В дальнейшем считаем $j\le n$, $a<0<b$. 


 


\begin{thm} Если число $q$ является четным, то система 


функций $\varphi_0,\dotsc,\varphi_n$ будет чебышевской системой, если 


$q$ нечетное, то любой ненулевой полином по этой системе имеет 


более $n+1$ нулей. 


\end{thm} 


 


\begin{pf} Для $P(x)=\sum\limits^n_{i=0}c_i\varphi_i(x)$ через $N^+(P)$ 


и $N^-(P)$ 


обозначим соответственно число его различных положительных 


и отрицательных корней, а через $W^+(P)$ ($W^-(P)$) --- число перемен 


знака полинома $P(t)$ ($P(-t)$), $N(P)$ --- число всех различных 


вещественных нулей $P$. Известно [10], что 


\begin{equation} 


N^+(P)\le W^+(P),\quad N(P)\le r+W^+(P)+W^-(P), 


\end{equation} 


где $r=1$, если $x=0$ является корнем $P$, и нулю в противном случае. 


 


Имеются три случая для полинома $P$. Первый, когда 


$c_0=\dotsb=c_{j-1}=0$, и второй случай, когда $c_j=\dotsb=c_n=0$, 


рассматриваются 


очевидным образом. Разберем третий случай, когда 


\begin{alignat*}{3} 


&c_0=\dotsb=c_{s-1}=0, &\quad 


c_s&\ne 0, &\ \; s&\ge 0,\\ 


%


&c_k=\dotsb=c_{j-1}=0, &\quad 


c_{k-1}&\ne 0, &\ \; k&\le j,\ \; k-1\ge s,\\ 


%


&c_j=\dotsb=c_{j+m-1}=0, &\quad  


c_{j+m}&\ne 0, &\ \; m&\ge 0, \\ 


%


&c_{n-i+1}=\dotsb=c_n=0, &\quad 


c_{n-j}&\ne 0, &\ \; i&\ge 0,\ \; n-i\ge j+m, 


\end{alignat*} 


причем в этих строках при $i=0$, $m=0$, $k=j$ или $s=0$ учитываются 


только неравенства с коэффициентами полинома $P$. В третьем 


случае полином имеет вид 


$$ 


P(x)=c_s x^s+\dotsb+c_{k-1} x^{k-1}+c_{j+m}x^{j+m+q} 


+\dotsb+c_{n-i}x^{n-i+q}. 


$$ 


 


Подсчитаем $W^+(P)$ и $W^-(P)$. Если полином 


$Q(x)=a_s x^s+\dotsb+a_n x^n$, $s\ge 0$, то 


\begin{equation} 


W^+(Q)+W^-(Q)\le n-s. 


\end{equation} 


При $s=0$ это неравенство приводится в ([10], 


гл.\,1); если же $s>0$, то $Q(x)=x^sQ_1(x)$, где 


$Q_1(x)=a_s+\dotsb+a_n x^{n-s}$. 


Разобьем последовательность коэффициентов полинома $P$ 


на части


$$ 


\Pi_1=(c_s,\dotsc,c_{k-1}),\quad 


\Pi_2=(c_{k-1},c_{j+m}),\quad 


\Pi_3=(c_{j+m},\dotsc,c_{n-i}). 


$$ 


Обозначим через $W^+(P,\Pi_i)$, $W^-(P,\Pi_i)$ число перемен знака $P$ 


на части $\Pi_i$, $i=1,2,3$. Тогда $W^+(P)+W^-(P)$ 


равен сумме получаемых 


таким образом шести чисел. 


 


Из неравенства (5) следует, что 


$$ 


W^+(P,\Pi_1)+W^-(P,\Pi_1)\le k-1-s,\quad 


W^+(P,\Pi_3)+W^-(P,\Pi_3)\le n-i-j-m. 


$$ 


Положим $W=W^+(P,\Pi_2)+W^-(P,\Pi_2)$, $R=j-k+1+m$, 


$a=c_{k-1}c_{j+m}$ и $b=(-1)^{k-1+j+m+q}c_{k-1}$, 


$c_{j+m}=(-1)^{R+q}a$. 


Очевидно, $W$ равен количеству отрицательных чисел в множестве 


$\{a,b\}$. 


Рассмотрим первую часть теоремы, когда $q$ --- четное число. 


Тогда $b=(-1)^R a$ и $W=1$ при $R$ нечетном; если же 


$R$ четное, то $W=0$ или $W=2$. 


 


Складывая все шесть чисел, получим 


$W^+(P)+W^-(P)\le n+W-(j-k+m+1)-s$. 


Число $r$ в неравенстве (4) не превосходит $s$. Значит, 


\begin{equation} 


W^+(P)+W^-(P)+r\le n+W-R. 


\end{equation} 


 


Непосредственно из обозначения $R$ видно, что $R\ge 1$, если 


оно нечетное, и $R\ge 2$ при $R$ четном. Поэтому из неравенств 


(4) и (6) получается первая часть теоремы. 


 


Вторая часть теоремы есть следствие ее первой части. 


Действительно, если бы какой-нибудь ненулевой многочлен по 


системе ($q$ нечетное) 


$$ 


1,\dotsc,x^{j-1},x^{j+q},\dotsc,x^{n+q} 


$$ 


имел бы более $n+1$ корней, то этот же многочлен, рассматриваемый 


по системе 


$$ 


1,\dotsc,x^{j-1},x^{j+q-1},\dotsc,x^{n+1+(q-1)},


$$ 


имел бы более $n+1$ корней, хотя по доказанной части она является 


чебышевской. 


\end{pf} 


 


\begin{note} Если $j=0$, то теорема остается в силе, 


когда $q$ нечетное; при $j=0$ и $q$ четном она неверна (например, 


система $x^2$, $x^3$, $x^4$ на отрезке $[-1,1]$ чебышевской не является). 


\end{note} 


 


\begin{note} Если из системы степеней выбрасывать не подряд идущие 


куски, то в оставшейся системе закономерности типа теоремы 4 


не соблюдаются. 


\end{note} 


 


{\bf 4.} Ниже все обозначения (в частности, 


$\varphi_0,\dotsc,\varphi_n$ п.\,3) сохраняются. 


 


\begin{thm} Пусть для пары $(g,G)$ множество $L(g,G)$ непустое. 


Тогда для любых $n\ge 2$, $q\ge 0$, $j\ge 1$ возможны только


следующие ситуации. 


\begin{itemize} 


\item[1)] $F(g,G)=1$, $q$ четное. Множество $L(g,G)$ содержит 


единственный полином, который является избыточным ужом пары 


$(g,G)$. 


\item[2)] $F(g,G)=1$, $q$ нечетное. Если $g_0$ имеет хотя бы два 


полинома наилучшего приближения по системе


$\varphi_0,\dotsc,\varphi_n$, то 


пара $(g,G)$ имеет бесконечно много ужей. Если же $g_0$ по этой 


системе имеет единственный полином наилучшего приближения, 


то он и будет единственным ужом пары $(g,G)$. 


\item[3)] $F(g,G)<1$. Тогда пара $(g,G)$ порождает ровно два ужа 


при $q$ четном и по меньшей мере два ужа для нечетного $q$. 


\end{itemize} 


\end{thm} 


 


\begin{pf} Для доказательства используются


теоремы 2, 3. Добавим только, что в случае, когда $q$ --- 


нечетное число, подсистема 


$1,\dotsc,x^{j-1},\dotsc,x^{j+q+1},\dotsc,x^{n-1+(q+1)}$ 


является чебышевской. 


\end{pf} 


 


Из следствия 1 теоремы 2 получается 


 


\begin{thm} Пусть 


$$ 


L=\bigg\{\,\sum^{j-1}_{i=0}c_i x^i+x^q 


\sum^n_{i=j}c_i x^i\bigg\} 


$$ 


--- множество полиномов по системе $\varphi_0,\dotsc,\varphi_n$, 


$F(g,G)<1$. Тогда 


решение задачи 


\begin{equation} 


c_k(l)\to\sup,\ \; l\in L(g,G), 


\tag{\text{з2}} 


\end{equation} 


является ужом, порожденным парой $(g,G)$ при $k=j-1$ $(j\ge 2)$ и 


$k=j+q$ $(j\ge 1)$, когда $q$ --- нечетное число. 


 


Если $q$ --- число четное, а $k=n+q$, то полином $l^*$ является 


решением 


задачи {\em (з2)} тогда и только тогда, когда $l^*$ является 


ужом с б\'ольшим старшим коэффициентом из двух ужей, порожденных


парой $(g,G)$. 


\end{thm} 


 


\begin{note} Задача (з2) при $q=0$ для $0\le k\le n$ изучалась 


ранее в [9]. 


\end{note} 


 


\begin{thebibliography}{99} 


 


\bibitem{1} 


Karlin S. {\em Representation theorems for positive functions} // 


J. Math. Mech. -- 1963. -- V.\,12. -- \No\,4. -- P.\,599--618. 


\bibitem{2} 


Дзядык В.К. {\em Введение в теорию равномерного приближения 


функций полиномами}. -- М.: Наука, 1975. -- 510 с. 


\bibitem{3} 


Долженко Е.П., Севастьянов Е.А. {\em Об определении 


чебышевских ужей} // Вестн. Моск. ун-та. Сер. матем., механ. -- 


1994. -- \No\,3. -- C.\,49--59. 


\bibitem{4} 


Долженко Е.П., Севастьянов Е.А. {\em Знакочувствительные 


аппроксимации $($пространство знакочувствительных весов, 


жесткость и свобода системы$)$} // Докл. РАН. -- 1993. -- 


Т.\,332. -- \No\,6. -- C.\,14--17. 


\bibitem{5} 


Долженко Е.П., Севастьянов Е.А. {\em Знакочувствительные 


аппроксимации $($вопросы единственности и устойчивости$)$} 


// Докл. РАН. -- 1993. -- T.\,333. -- \No\,1. -- С.\,5--7. 


\bibitem{6} 


Долженко Е.П., Севастьянов Е.А. {\em Метрические пространства 


полунепрерывных функций} // Матем. заметки. -- 1994. -- T.\,55. -- 


\No\,3. -- C.\,48--59. 


\bibitem{7} 


Загиров Н.Ш. {\em Рациональные ужи} // Матем. заметки. -- 1993. 


-- T.\,53. -- \No\,6. -- C.\,33--40. 


\bibitem{8} 


Загиров Н.Ш. {\em Несимметричная норма} // Сб. ``Kонструктивная 


теория функций и ее приложения''. -- Махачкала, 1994. -- C.\,47--49. 


\bibitem{9} 


Zagirov N.Sh. {\em Polynomials with extremal coefficients and 


snakes} // East J. on Approximations. -- 1995. -- V.\,1. -- 


\No\,2. -- P.\,231--248. 


\bibitem{10} 


Полиа Г., Сег\"е Г. {\em Задачи и теоремы из анализа}. Ч.\,2. -- 


3-е изд. -- M.: Наука, 1978. -- 432 с. 


 


\end{thebibliography} 


 


\vskip 2cm 


 


\postup{\it Дагестанский государственный}{$\qquad$\it Поступили} 


\postup{\it университет}{{\it первый вариант} 08.07.1996} 


\postup{}{{\it окончательный вариант} 22.04.1997} 


 


\label{end} 


\end{document} 


